\chapter{Empirical Attribution of Failure Modes}
\label{ch:empirical}



In the previous chapter, we organised MP-VRDU by how systems manage evidence they cannot hold all at once. The field has explored this problem from many directions, and its findings have not settled into agreement. 

A large body of work is OCR-free (Table~\ref{tab:model_survey_arch}), rendering pages or page regions and reading them with vision on the premise that image input sidesteps the cascading errors of text extraction and preserves layout information that parsing discards \citep{tanaka2025vdocrag, zheng2026doc, cho2024m3docrag}. Against this sits measured counter-evidence that a reasoner given no extracted text hallucinates on dense passages rather than reading them \citep{hannan2026docslm}, and that combining both channels beats either alone \citep{han2025mdocagent, suri2025visdom}. Yet other work measures a vision-only pipeline outperforming exactly such a combination, leaving it unclear whether text earns its cost \citep{liu2025resolving}.

How much evidence to place before the reasoner is contested in turn. Fixed-depth retrieval is widely reported as brittle in both directions. Retrieving too few units omits required evidence and caps accuracy regardless of the reasoner \citep{yan2026docseeker, li2026avir}. Retrieving too many admits distractors that dilute the reasoner's attention, and evidence at some positions in a long context is used less reliably than evidence at others \citep{liu2024lost, wang2024needle}, so answers degrade even when the required pages are present \citep{pez2025enhancing, yan2026docseeker}. Which failure dominates is unresolved: some work identifies retrieval coverage as the single factor that most improves accuracy \citep{liu2025resolving, zhu2026doclens}, while other work treats dilution as serious enough to require reranking or adaptive depth rather than a larger fixed pass \citep{li2026avir, pez2025enhancing}.

What the MLLM reasoner does with the evidence it holds is the most broadly contested of all. One line of work places the failure in cross-page assembly, where a model that reads each page correctly still fails to combine them \citep{zheng2026doc}. Another places it within a single page, in a shallow comprehension that locates the right evidence and then misreads a chart or miscalculates over a table \citep{guo2026end}. A separate dispute concerns faithfulness. Models hallucinate, answering questions the evidence cannot support rather than abstaining \citep{landeghem2023document, zhang2025dream}, and the field agrees neither on whether answering or abstaining is the safer default nor on how to control the choice \citep{ma2024mmlongbench}.

Such claims tend to enter a system as a design premise rather than a tested result. They often point in opposite directions, and they are reported over pipelines, datasets, and operating points too different to compare directly. The field has gathered claims about how to build these systems faster than it has established which of them hold. The three disputes are separable because each concerns a different operation: encoding the document, selecting the evidence placed before the reasoner, and reasoning over that evidence. We take these as three loci of failure: representation, selection, and reasoning (\S\ref{sec:attribution_framework}), and in this chapter, we test the competing claims at each locus by isolated each in a single-pass pipeline, intervening on one while holding the others fixed (\S\ref{sec:emp_methodology}, \S\ref{sec:emp_findings}). We further conduct error analysis (\S\ref{sec:emp_error_analysis}), and we close by  discussing practical implications of designing and deploying MP-VRDU systems under a fixed compute budget (\S\ref{sec:practical_discussion}).

% Representation findings:
% The converse is equally important. Although vision expands the range of evidence available to the system, the page image alone remains below the combined representation, reaching 45.6\% against 52.5\% for TLV overall. This gap does not necessarily imply that the visual representation itself lacks the information required to answer these questions. In principle, much of the textual and structural evidence available through TLV is also visible in the page image, but recovering and using it may introduce additional limitations. Resolution may fail to preserve sufficiently fine-grained details, as examined further in our conversion-fidelity experiments; relevant evidence may occupy only a small region of an otherwise information-dense page, requiring finer-grained selection that introduces a separate retrieval problem; and the limitation may also lie in the vision-language model's ability to recognise, interpret, or reason over visual evidence. Text may therefore remain complementary not only because it provides different information, but because it presents information in a form that is currently easier for the system to recover and use. Accordingly, TLV's advantage should be interpreted as evidence of practical complementarity between the two representations under the evaluated pipeline, rather than demonstrating that vision and text contain inherently separate information.